% =============================================================================
% Revised results sections for Zito & Kowal, "A dynamic copula model for
% probabilistic forecasting of non-Gaussian multivariate time series."
% Ethan Carlson. Rerun of 20 November 2026, seed 8675309, k = n, 5000/5000.
%
% Figures come from figures/*.pdf (export_paper_figures.R) and table bodies
% from tables/*.tex (export_comparison_tables.R).
%
% PATCH MANIFEST (in reading order)
%   1. Sec. 3.4 (p. 11)   INSERT 1 sentence after MGP sentence; REPLACE
%                         k = \lceil 0.7n \rceil with k = n.
%   2. Sec. 5   (p. 19)   INSERT 3 paragraphs at end of section (SV-BVAR and
%                         explicit BVAR lag orders).
%   3. Sec. 6   (p. 20)   REPLACE item 2 of the two-item preamble list.
%   4. Sec. 6.1 (p. 21)   INSERT 2 sentences on T = 200 mixing at end of para 2.
%   5. Sec. 6.2 (p. 22)   REPLACE entire subsection with inferred-vs-fixed
%                         margins study (Figure 5, Table 1). Old Fig. 5 deleted.
%   6. Sec. 7.1 (p. 24)   INSERT 1 paragraph after Fig. 6 discussion (ablation).
%   7. Sec. 7.1 (p. 25)   REPLACE forecast metrics paragraph + Table 2.
%   8. Sec. 7.2 (p. 27)   INSERT 1 paragraph after Fig. 7 discussion (ablation),
%                         with Figure 8 and Table 5.
%   9. Sec. 7.2 (p. 28)   REPLACE forecast metrics paragraph with 4 paragraphs,
%                         Figure 9, Table 3, and Table 4.
%
% MASTER / SUPPLEMENT HOUSEKEEPING ON MERGE
%   - Sec. 7.2 setup (p. 27): add SV-BVAR alongside BVAR and DLM.
%   - Fig. 7 caption (p. 28): update "Abbreviations are in Table 4" to the
%     renumbered Appendix E variable table.
%   - Supplement Appendix C (p. 9): update k = \lceil 0.7n \rceil -> k = n.
%   - Renumber existing supplement figs. 8-10 and tables 3-4 to accommodate new 
%     main-text Figs. 8-9 and Tables 4-5.
%
% Strip the preamble helper macros and everything below the STRIP marker before
% pasting into the master manuscript.
% =============================================================================

\documentclass[11pt]{article}
\nofiles % no .aux, .toc, .lof, .lot. Only .log and .pdf.
\usepackage[margin=1in]{geometry}
\usepackage{booktabs}
\usepackage{amsmath}
\usepackage{xcolor}

% graphicx pulls epstopdf-base, which wants grfext.sty, missing in this TinyTeX.
% All figures are PDFs, so mark it loaded and skip it.
\expandafter\def\csname ver@epstopdf-base.sty\endcsname{2020/01/24}
\usepackage{graphicx}
\graphicspath{{figures/}}

\renewcommand{\topfraction}{0.9}
\renewcommand{\bottomfraction}{0.8}
\renewcommand{\textfraction}{0.07}
\renewcommand{\floatpagefraction}{0.7}
\newcommand{\todo}[1]{\textcolor{red}{[#1]}}

% --- STRIP PREAMBLE HELPERS ON PASTE -----------------------------------------
% Dummy labels and natbib-compatible lookup macros so \ref, \citet, and \citep
% render with exact full-paper numbers in standalone mode under \nofiles.
\makeatletter
\newcommand{\dummyref}[2]{\@namedef{r@#1}{{#2}{1}}}
\newcommand{\defcite}[3]{\@namedef{citet@#1}{#2}\@namedef{citep@#1}{#3}}
\newcommand{\citet}[1]{\@ifundefined{citet@#1}{[#1]}{\@nameuse{citet@#1}}}
\newcommand{\citep}[1]{\@ifundefined{citep@#1}{(#1)}{\@nameuse{citep@#1}}}
\makeatother

\dummyref{sec:priors}{3.4}
\dummyref{sec:sim_margins}{6.2}
\dummyref{sec:crime}{7.1}
\dummyref{fig:margin_cov}{5}
\dummyref{fig:ablation}{8}
\dummyref{fig:sharpness}{9}
\dummyref{tab:margin_cov}{1}
\dummyref{tab:crime}{2}
\dummyref{tab:macro}{3}
\dummyref{tab:dm}{4}
\dummyref{tab:ablation}{5}

\defcite{diebold1995}{Diebold and Mariano (1995)}{(Diebold and Mariano, 1995)}
\defcite{gneiting2007}{Gneiting et al. (2007)}{(Gneiting et al., 2007)}
\defcite{gruber2025}{Gruber and Kastner (2025)}{(Gruber and Kastner, 2025)}
\defcite{bayesianVARs}{Gruber et al. (2026)}{(Gruber et al., 2026)}
% -----------------------------------------------------------------------------

% grey is the manuscript as it stands, black is new. Locating only, not for
% pasting.
\definecolor{papergray}{gray}{0.42}
\newcommand{\paper}[1]{{\color{papergray}\textit{#1}}}
\newenvironment{splice}[1]{%
  \par\medskip\noindent{\small\bfseries #1}\par\smallskip
  \begingroup\leftskip=1.2em\relax\parindent=0pt\relax\sloppy}%
  {\par\endgroup\medskip}

\newcommand{\missing}[1]{\textcolor{red}{[\texttt{\detokenize{#1}} missing: run
  \texttt{\detokenize{export_paper_figures.R}}]}}
\newcommand{\intab}[1]{\IfFileExists{tables/#1.tex}{\input{tables/#1}}{\missing{tables/#1.tex}}}
\newcommand{\infig}[1]{\IfFileExists{figures/#1.pdf}{\includegraphics[width=\textwidth]{#1.pdf}}{\missing{figures/#1.pdf}}}

\begin{document}

% =============================================================================
\section*{Section 3.4, priors and computation}

\begin{splice}{INSERT, one sentence, after the multiplicative gamma process
sentence (p.~11).}

\paper{For the factor loadings, we use the multiplicative gamma process prior
(Bhattacharya and Dunson, 2011), which provides ordered shrinkage across the
columns of $\boldsymbol{\Lambda}$ to reduce sensitivity to the specified number
of factors $k$.}

We therefore fix $k$ at its saturated value $k = n$ and rely on the prior to
shrink the redundant loadings toward zero, which removes the choice of factor
dimension altogether rather than deferring it to a rule of thumb;
Appendix~\todo{ref} reports forecasts across a range of $k$.

\paper{The variances are assigned $v_i^{-1} \sim$ iid Gamma$(\alpha_0,
\beta_0)$. Next, we use a conditionally conjugate prior for the VAR parameters
in (6) \dots}

\end{splice}

\begin{splice}{REPLACE, two characters, four sentences later in the same
paragraph (p.~11).}

\paper{In each of the simulations and applications that follow, we use the same
default prior hyperparameters:} $k = n$\paper{, $d_0 = k + 1$,
$\boldsymbol{\Psi}_0 = \mathbf{I}_k$, $\mathbf{G}_0 = \mathbf{0}$,
$\mathbf{O}_0 = \mathbf{I}_k$, $\alpha_0 = 1$, $\beta_0 = 0.3$. Further details
are in the supplement.}

\end{splice}

% The simulations have n = 2, where ceil(0.7n) already equals n, so the "same
% defaults everywhere" claims in 6.1 and the conclusion are unaffected.

% =============================================================================
\setcounter{section}{4}
\section{Competing methods}

\begin{splice}{INSERT, three paragraphs, at the end of Section 5 (p.~19).}

\paper{By default, the static parameters $\mathbf{W}$ (when included) and
$\mathbf{V}$ are estimated via maximum likelihood.}

\medskip
Finally, for the macroeconomic application we add a Bayesian VAR with
stochastic volatility (\textbf{SV-BVAR}), which generalizes~(13) by allowing the
innovation variances to evolve over time. Time-varying volatility is a central
feature of macroeconomic data and a standard ingredient of modern forecasting
VARs, which makes this the most demanding of our Gaussian comparisons. We use
the implementation of \citet{gruber2025} in the \texttt{bayesianVARs}
package \citep{bayesianVARs} in \texttt{R}, with $p = 4$ lags, a horseshoe
prior on the autoregressive coefficients, a factor stochastic volatility
specification for the innovation covariance, and 5{,}000 draws retained after
5{,}000 burn-in.

This specification places a latent factor structure in the innovation
covariance, where the DGFC places one in the copula, meaning the two models are
factor models of different objects. Posterior draws are restricted to those
satisfying the stationarity condition, which mirrors the truncation the DGFC
imposes on its own factor VAR in~(6) and makes the two models comparable on
that dimension.

The BVAR of~(13) uses $p = 4$ for the macroeconomic data and $p = 1$ for the
crime counts, in both cases with the default hierarchical Minnesota prior of
the \texttt{BVAR} package.

\end{splice}

% =============================================================================
\section{Simulation results}

\begin{splice}{REPLACE item 2 of the two-item list, end of the Section 6
preamble (p.~20). Item 1 and the lead-in are unchanged.}

\paper{These DGPs present significant challenges: in each case, the DGFC has
misspecified dynamics and no knowledge about the marginal stationary
distributions. In that context, we seek to answer two main questions:}

\begin{enumerate}
\item \paper{Does the margin adjustment posterior provide accurate estimates
with precise and well-calibrated uncertainty quantification for unknown marginal
stationary distributions?}
\item Does inferring the margins improve probabilistic forecasts relative to the
common practice of fixing them at a point estimate?
\end{enumerate}

\end{splice}

% Item 2 was: "Does our (pseudo-)posterior predictive distribution provide
% accurate probabilistic forecasts compared to popular MTS alternatives?"
% Swapped, not added: that comparison now lives entirely in Section 7.

\subsection{Recovering the marginal stationary distributions}

\begin{splice}{INSERT, two sentences, at the end of the second paragraph of 6.1
(p.~21). Note the grey sentence it has to sit beside.}

\paper{Notably, convergence is typically quite fast: $T = 50$ or $T = 100$ is
enough to achieve strong concentration around the ground truth. Taken together,
these results demonstrate the ability of the margin adjustment posterior to
quickly and automatically learn the marginal stationary distributions of MTS
exhibiting a range of distributional features (discrete support, heavy tail,
skew, etc.) and across a wide variety of underlying dynamics.}

The posterior mean is unaffected, but one caveat applies to the bands at the
largest sample size. The truncation intervals in the latent sampling step narrow
as the observations fill in, so the MCMC budget required for the margin
posterior grows with $T$, and the bands at $T = 200$ are consequently somewhat
tighter than the pseudo-posterior they approximate (Appendix~\todo{ref}).

\end{splice}

% "convergence is typically quite fast" and "quickly and automatically learn"
% both need a hedge, or the caveat undercuts Figure 4 two sentences after the
% paper finishes praising it. Dan's half; flag rather than edit.

\subsection{Inferred versus fixed margins}

% Replaces the old 6.2 outright. Figure 5 and its discussion are scrapped.

Common practice in the MTS copula literature is to fix the margins at a point
estimate and proceed as though they were known, where the margin adjustment
infers $F_{1:n}$ jointly with the copula parameters. To isolate what that extra
layer of uncertainty buys, we rerun the forecasting exercise twice from the same
posterior fit---once with the margin adjustment, and once with $F_{1:n}$ frozen
at the empirical CDF of the training window---holding everything else fixed.

We use the same three DGPs, sample sizes $T = 25, 50, 100, 200$, and 100
replicates per cell with the DGP parameters redrawn on each replicate. Both arms
are forecast from a single posterior fit with the random number state restored
between calls, so that the latent forward simulations are identical and the only
difference between the two is the treatment of the margin. The comparison is
therefore paired at the level of the individual forecast, and because each DGP
has a closed-form stationary distribution we can also score each estimated
margin against the truth.

% \caption still runs off LaTeX's own counter, need to set to match the full-paper
% slot left free by the scrapped Figure 5.
\setcounter{figure}{4}
\begin{figure}[!htbp]
\centering
\infig{fig_margin_coverage}
\caption{Coverage of the 95\% forecast interval at $h = 1$ against training
sample size, for the margin adjustment and for a margin frozen at the empirical
CDF. Each point averages 200 forecasts (100 replicates, two variables). The
dashed line is nominal coverage.}
\label{fig:margin_cov}
\end{figure}

\begin{table}[!htbp]
\centering\small
\intab{tab_margin_coverage}
\caption{Forecast interval coverage at $h = 1$, pooled across the three DGPs
and both variables, 600 forecasts per cell. Shortfall is the distance below the
0.95 nominal level. The $p$-value is from an exact McNemar test on the paired
coverage indicators, which is available because both arms are scored on the
same forecast origins from the same posterior fit.}
\label{tab:margin_cov}
\end{table}

The benefit shown in Figure~\ref{fig:margin_cov} and Table~\ref{tab:margin_cov}
is concentrated in small samples and is specific to interval calibration.
Inferring the margins closes roughly three fifths of the coverage shortfall at
$T = 25$ and removes it entirely by $T = 50$, and an exact McNemar test on the
paired coverage indicators still separates the two arms at $T = 100$. By
$T = 200$ they are indistinguishable.

The gain does not extend to the other metrics. Measured by mean absolute error
against the true stationary CDF at the observed values, the inferred margin is a
four to fourteen percent worse estimate of $F$ itself from $T = 50$ onward in
all three DGPs. Density forecasts divide by DGP rather than by sample size,
freezing $F$ being consistently better under the VARCH and consistently worse
under the VARMA by a few percent either way, and neither under the VARMA copula,
with no trend in $T$ in any of the three. The benefit of inferring the margins
is therefore confined to interval calibration.

A likely mechanism is that an empirical CDF assigns no probability outside the
observed range, and that realized values fall outside that range less often as
$T$ grows. The simulation offers only weak support for this---the correlation
between the exceedance rate and the CRPS penalty is 0.24---because none of the
three DGPs is persistent enough to leave its observed range often. The
applications of Section~7 are more informative. The shortest macroeconomic
training window is longer than the shortest crime window, yet the macroeconomic
ablation shows the larger loss of coverage, which points to the type of series
rather than to its length.

% =============================================================================
\setcounter{section}{6}
\section{Real-data applications}

\subsection{Crime counts in New South Wales, Australia}

\begin{splice}{INSERT, new paragraph, after the Figure 6 discussion (p.~24--25).
It answers the grey sentence marked with a dagger, which currently rests on
eyeballing the left panel.}

\paper{Figure 6 displays summaries of the parameter inferences from our model.
On the left we plot the margin adjustment posterior for a selection of variables
that exhibit heterogeneous features (small counts, heavy tails, asymmetry).
$^\dagger$Using the empirical cumulative distribution function (ECDF) as a
reference point, the margin adjustment appears adequate while also providing
reasonable uncertainty quantification. On the right, we provide heat map
posterior summaries \dots In general, cross-sectional correlations persist
serially, while autocorrelations are strong for each individual series.}

\medskip
Freezing the margins at the empirical CDF narrows the 95\% intervals by five
percent on average while lowering empirical coverage from 0.958 to 0.944, so the
sharper intervals are bought by under-covering rather than by better
calibration. The rerun holds the posterior fit and everything else fixed, which
isolates the contribution of inferring $F_{1:n}$ and puts a number on the visual
comparison in the left panel of Figure 6 (Section~\ref{sec:sim_margins} does
the same on simulated data). The loss is uneven. Three series fall well below
nominal once the margin is frozen---blackmail, escaping custody and resisting
an officer---while non-domestic assault, the highest-count series, is the one
variable on which the frozen margin covers slightly better
(Figure~\ref{fig:ablation}, left, and Table~\ref{tab:ablation}). The series
that are unaffected are those the margin adjustment already over-covers,
accessory to murder and manslaughter among them, where there is room to absorb a
narrower interval without falling short of nominal. Density forecasts move in
the same direction as coverage, with eight of the ten series preferring the
inferred margin.

\end{splice}

\begin{splice}{REPLACES the whole paragraph on p.~25--26. The original is the
grey block.}

\paper{Table 1 presents one-month-ahead forecast metrics for each variable,
averaged over the period January 2000 ($t_0$) to December 2023 ($T$). We compare
our approach to BVAR($p = 1$) and two count-valued multivariate state space
models: a Poisson DLGM (P-DGLM) and a negative binomial DGLM (NB-DGLM), both
using the dynamics from (15). For all variables, the DGFC delivers the best
density forecasts and the smallest 95\% interval forecasts with near nominal
coverage. For nine of ten variables, the DGFC provides superior point forecasts,
often by a large amount.}

\medskip
% "Poisson DLGM" is a typo in the original, corrected below.
Table~\ref{tab:crime} presents one-month-ahead forecast metrics for each
variable, averaged over the period January 2000 ($t_0$) to December 2023 ($T$).
We compare our approach to BVAR($p=1$) and two count-valued multivariate state
space models: a Poisson DGLM (P-DGLM) and a negative binomial DGLM (NB-DGLM),
both using the dynamics from~(15). For all variables, the DGFC delivers the
best density forecasts and the smallest 95\% interval forecasts with near
nominal coverage. For nine of ten variables, the DGFC provides superior point
forecasts, often by a large margin; blackmail is the lone exception, where the
BVAR attains a smaller mean absolute error at a far greater interval width. The
NB-DGLM is the only competitor whose coverage is comparable, and it purchases
that coverage with intervals wider than the DGFC's by factors ranging from
three to nearly fifty. Diebold--Mariano tests on the per-origin CRPS
differentials \citep{diebold1995}, computed with Newey--West standard errors,
confirm that these gaps are not an artifact of sampling noise
(Table~\ref{tab:dm}). Our approach is significantly better on all ten series
against the BVAR and the NB-DGLM, on nine of ten against the P-DGLM, and
significantly worse on none.

\end{splice}

\begin{table}[!htbp]
\centering\small
\setlength{\tabcolsep}{4pt}
\intab{tab_crime}
\caption{One-month-ahead forecast metrics for the crime count data in
Figure~1, averaged over the period January 2000 to December 2023. \textbf{MAE}
is the mean absolute error of the posterior predictive median (smaller is
better). \textbf{COV} and \textbf{SIZE} are the empirical coverage and average
length of 95\% credible intervals based on posterior predictive quantiles
(prefer small intervals with high coverage). \textbf{CRPS} is the continuous
ranked probability score of the entire forecast distribution (smaller is
better). The best (smallest) values of MAE, SIZE, and CRPS are bolded for each
variable. A CRPS set in italic and marked $^{\dagger}$ is one that a 
Diebold--Mariano test on the per-origin CRPS differentials, with Newey--West 
standard errors, cannot separate from the DGFC at the 5\% level.}
\label{tab:crime}
\end{table}

\clearpage

% =============================================================================
\subsection{US macroeconomic aggregates}

\begin{splice}{INSERT, new paragraph, after the Figure 7 discussion (p.~27--28).}

\paper{Figure 7 displays summaries of the parameter inferences in our model.
Again, the DGFC captures the non-Gaussian marginals and provides reasonable
uncertainty quantification. The estimated latent correlations identify high
positive associations the real economic activity variables (output,
productivity, etc) \dots These stylized results match various basic economic
intuitions about the dependence structure of these variables.}

\medskip
Repeating the ablation of Section~\ref{sec:crime} on the macroeconomic data
(Figure~\ref{fig:ablation}, right, and Table~\ref{tab:ablation}) yields a
stronger version of the same result. Freezing the margins at the empirical CDF
narrows the intervals on eleven of the fourteen series and drops mean coverage
from 0.946 to 0.914, more than twice the deterioration observed on the crime
counts. Density forecasts divide by variable type. On eight of the ten growth
rates the frozen margin is the better of the two, by as much as nine percent,
while on capacity utilization, the unemployment rate and the federal funds rate
it is worse by 37, 99 and 69 percent respectively. The three series on which it
fails badly are levels---bounded and highly persistent, and therefore precisely
the series for which the observed range over a training window is least
informative about the support of the predictive distribution. Inferring the
margins is thus not uniformly better on these data. It costs little where the
empirical CDF is adequate and helps substantially where it is not.

\end{splice}

% "high positive associations the real economic activity variables" in the grey
% block is missing a word. Worth fixing while that page is open.

% Figure 8 follows existing Figure 7 and precedes Figure 9 in citation order.
\setcounter{figure}{7}
\begin{figure}[!htbp]
\centering
\infig{fig_ablation}
\caption{Empirical coverage of the 95\% forecast intervals under the margin
adjustment and under a margin frozen at the empirical CDF, for the crime counts
(left) and the macroeconomic aggregates (right). Both arms come from the same
posterior fit at each forecast origin, so the only difference is the treatment
of the margin. The dashed line is nominal coverage.}
\label{fig:ablation}
\end{figure}

% Set counter to 4 so standalone caption renders as Table 5; in the master
% manuscript, place this float at the end of Section 7.2 after Tables 3 and 4.
\setcounter{table}{4}
\begin{table}[!htbp]
\centering\small
\intab{tab_ablation}
\caption{The empirical CDF ablation, pooled over series within each
application. Width is the mean change in 95\% interval length when the margin
is frozen, so a negative figure means the frozen margin is sharper.}
\label{tab:ablation}
\end{table}

\begin{splice}{REPLACES the whole paragraph on p.~28, expanding one paragraph
into four.}

\paper{Table 2 displays the one-step ahead forecasting results, averaged over
the period 1985Q1 ($t_0$) to 2014Q1 ($T$). Across all variables, the DGFC
delivers the smallest forecasting intervals---often by a wide margin---with
close to nominal (95\%) coverage. By contrast, the DLM and BVAR are excessively
conservative. Similarly favorable results are observed for density forecasting,
with CRPS preferring DGFC for the majority of variables. Point forecasting
performance tends to favor the DLM, although the DGFC remains competitive. These
results suggest that the DGFC may be most useful for interval and density
forecasting, both of which are essential for probabilistic uncertainty
quantification.}

\medskip
% Three claims in the grey block no longer hold. "Across all variables" is now
% eleven of fourteen on SIZE. "Point forecasting tends to favor the DLM" is
% reversed. And the closing sentence is the one the SV-BVAR tie complicates.
Table~\ref{tab:macro} displays the one-step-ahead forecasting results, averaged
over the period 1985Q1 ($t_0$) to 2014Q1 ($T$), and Figure~\ref{fig:sharpness}
summarizes the interval forecasts according to the principle of maximizing
sharpness subject to calibration \citep{gneiting2007}. Against the BVAR and the
DLM, the picture is much the same as for the crime counts. The DGFC delivers
the smallest interval forecasts---often by a wide margin---with close to
nominal coverage, while the BVAR covers at unity on every series and the DLM
remains excessively conservative. Density forecasts favor the DGFC over the
BVAR on all fourteen series and over the DLM on twelve, with neither competitor
significantly better on any. Point forecasting, which in earlier versions of
this exercise tended to favor the DLM, now favors the DGFC on eight of the
fourteen series, a reversal consistent with saturating the factor dimension
(Section~\ref{sec:priors}) and concentrated in the level variables.

The SV-BVAR is the more demanding comparison, and against it the result is a
draw rather than a win. Our approach produces narrower intervals on eleven of
the fourteen series, with mean coverage nearer the nominal level, and point
accuracy divides evenly between the two. On density forecasts the two methods
are statistically indistinguishable on twelve of the fourteen series by a
Diebold--Mariano test on the per-vintage CRPS differentials
(Table~\ref{tab:dm}); the SV-BVAR is significantly better on the remaining two,
and our approach on none. That the DGFC should hold its own against a model
with time-varying volatility---a prominent feature of macroeconomic data, and
one the DGFC makes no attempt to capture---suggests that much of what it gains
in this exercise comes from the non-Gaussian margins and the latent dependence
structure rather than from conditional heteroskedasticity.

The two series on which the SV-BVAR prevails, the yield spread and the federal
funds rate, are levels rather than growth rates, and the pattern extends to the
remaining level variables. The SV-BVAR's CRPS advantage is positive for all four
levels in every subperiod of the sample, while on the ten growth rates the two
methods trade places. Nor is the advantage driven by volatile episodes, being
smallest over 2007--2009.

A plausible explanation is that these are the most persistent series in the
cross-section, and that our model is stationary by construction where the
SV-BVAR is not. Discarding the SV-BVAR's explosive autoregressive draws does not
render that model stationary, since its innovation variances remain free to
wander, whereas the DGFC enjoys no comparable latitude. The restriction is
visibly binding on both sides---the SV-BVAR retains a median of 2{,}828 of its
5{,}000 draws, and fewer than half of them at 47 of the 116 vintages, while the
DGFC's factor VAR sits at a mean spectral radius of 0.977 across the
posterior---but only one of the two models is obliged to remain inside it. Lag
order is a secondary difference, in that~(6) restricts the latent factor VAR to
a single lag. We leave a more flexible specification to future work; as matters
stand, the cost of the restriction is confined to two of fourteen series.

\end{splice}

\begin{figure}[!htbp]
\centering
\infig{fig_sharpness}
\caption{Sharpness against calibration for the fourteen macroeconomic series.
Each point is one series under one model. The horizontal position is the average
95\% interval width relative to the DGFC's width for that series, on a log
scale, and the vertical position is empirical coverage. Sharper is to the left
and calibrated is on the dashed line, so the preferred region is to the left
along that line. Widths are normalized because the series differ in scale by
three orders of magnitude.}
\label{fig:sharpness}
\end{figure}

\setcounter{table}{2}
\begin{table}[!htbp]
\centering\small
\setlength{\tabcolsep}{4pt}
\intab{tab_macro}
\caption{One-quarter-ahead forecast metrics for the macroeconomic data in
Figure~2, averaged over the period 1985Q1 to 2014Q1. \textbf{MSE} is the mean
squared error of the posterior predictive mean (smaller is better).
\textbf{COV}, \textbf{SIZE}, and \textbf{CRPS} are defined in
Table~\ref{tab:crime}. The best (smallest) values of MSE, SIZE, and CRPS are 
bolded for each variable. A CRPS set in italic and marked $^{\dagger}$ is one 
that a Diebold--Mariano test on the per-vintage CRPS differentials, with 
Newey--West standard errors, cannot separate from the DGFC at the 5\% level.}
\label{tab:macro}
\end{table}

% \begin{table}[!htbp]
% \centering\small
% \intab{tab_dm}
% \caption{Diebold--Mariano tests on the CRPS differential against each
% competitor, computed with Newey--West standard errors. Negative means the DGFC
% is better. Entries are significant at the 5\% level unless set in italic and marked
% $^{\dagger}$, so the daggers locate the comparisons the test cannot separate.
% The macroeconomic differentials are per vintage, the crime differentials per
% forecast origin.}
% \label{tab:dm}
% \end{table}

\clearpage

% =============================================================================
% STRIP BELOW THIS LINE
% =============================================================================

\section*{One open decision}

Whether the 6.1 mixing caveat sits in the body or the appendix. Professor Kowal 
noted not to over-advertise MCMC efficiency cuts in favor of stating it somewhere. 
The short version is in 6.1 above, while the longer version, which says what the 
mixing costs, is below. Let me know what you both think of my edits and insertions.

\begin{quote}
One caveat applies to the credible bands at the largest sample size. The
truncation intervals in the latent sampling step narrow as the observations fill
in, so the MCMC budget required for the margin posterior grows with $T$ while
the budget required for the predictive does not. At the default 10{,}000 sweeps
the band at $T = 200$ covers the true stationary CDF between 79 and 93 percent
of the time depending on the DGP, rising to between 86 and 96 percent at
160{,}000 sweeps and still widening. The posterior mean is unaffected, and so
are the forecasts, whose interval coverage at $T = 200$ runs between 0.925 and
0.975 on the same fits. The default is adequate for every result reported here;
credible bands for $F$ at $T = 200$ are not among them.
\end{quote}

\section*{New BibTeX entries}

\begin{verbatim}
@article{diebold1995,
  author  = {Diebold, Francis X. and Mariano, Roberto S.},
  title   = {Comparing Predictive Accuracy},
  journal = {Journal of Business \& Economic Statistics},
  year    = {1995},
  volume  = {13},
  number  = {3},
  pages   = {253--263}
}

@article{gneiting2007,
  author  = {Gneiting, Tilmann and Balabdaoui, Fadoua and Raftery, Adrian E.},
  title   = {Probabilistic Forecasts, Calibration and Sharpness},
  journal = {Journal of the Royal Statistical Society: Series B},
  year    = {2007},
  volume  = {69},
  number  = {2},
  pages   = {243--268},
  doi     = {10.1111/j.1467-9868.2007.00587.x}
}

@article{gruber2025,
  author  = {Gruber, Luis and Kastner, Gregor},
  title   = {Forecasting Macroeconomic Data with Bayesian {VAR}s: Sparse or Dense? {I}t Depends!},
  journal = {International Journal of Forecasting},
  year    = {2025},
  volume  = {41},
  number  = {4},
  pages   = {1589--1619},
  doi     = {10.1016/j.ijforecast.2025.02.001}
}

@manual{bayesianVARs,
  author = {Gruber, Luis and Haan, Stefan and Kastner, Gregor},
  title  = {bayesianVARs: {MCMC} Estimation of Bayesian Vectorautoregressions},
  year   = {2026},
  note   = {R package},
  doi    = {10.32614/CRAN.package.bayesianVARs}
}
\end{verbatim}

% gneiting2007 was the JASA scoring-rules paper, which is a different article.
% The sentence citing it is about maximizing sharpness subject to calibration,
% which is the JRSS-B paper above. If CRPS is cited anywhere else on merge,
% that one does want Gneiting and Raftery (2007), JASA 102(477), 359-378.

\end{document}